\documentclass[11pt,letterpaper]{amsart}
\usepackage[T1]{fontenc}
\usepackage{lmodern}
\usepackage[hmargin=1.15in,vmargin=0.95in]{geometry}
\usepackage{microtype}
\usepackage{booktabs,tabularx}
\usepackage[hidelinks]{hyperref}
\hypersetup{
  pdftitle={An existence classification for eight-vertex origami tori},
  pdfauthor={The-Anh Vu-Le}
}

\newtheorem{theorem}{Theorem}
\newcommand{\R}{\mathbb R}
\newcommand{\Z}{\mathbb Z}
\newcommand{\C}{\mathbb C}
\newcommand{\covol}{\operatorname{covol}}
\newcommand{\dd}{\mathrm d}
\newcommand{\release}{b391bdb86a7287c6034a77b1b3792cabd00425bb}
\newcommand{\src}[2]{\href{https://github.com/vltanh/lean4-origami-torus/blob/\release/#1}{\texttt{#2}}}

\title[Eight-vertex origami tori]{An existence classification for eight-vertex origami tori}
\author{The-Anh Vu-Le}
\date{Draft, November 27, 2026}
\subjclass[2020]{52C25, 57Q05}
\keywords{Paper torus, flat torus, triangulation, determinant signs, formal proof}

\begin{document}

\begin{abstract}
Of the seven simplicial triangulations of the torus on eight vertices,
exactly Types \(3\)--\(7\) in Lutz numbering admit embedded realizations
in \(\R^3\) with all cone angles equal to \(2\pi\). The remaining
existence question concerned Type \(2\). We reduce a hypothetical flat
embedding to a finite system of determinant signs and give a
kernel-checked propositional refutation. A diagonal exchange in the
planar families of Doyle, Lander, Schwartz, and Trettel, followed by
inflation at fixed periods, gives a universal construction on Type \(6\).
Types \(4\) and \(7\) realize every reduced modulus off the sharp-rhombic
ray. For Type \(5\) we give exact radical coordinates for the square
torus and a uniform radical construction on a closed box of
nonrectangular moduli. The nonexistence and realization theorems are
formalized in Lean~4.
\end{abstract}

\maketitle

\section{Introduction}

A flat torus is a quotient \(T_\Lambda=\C/\Lambda\) by a rank-two
lattice, with the metric induced by the Euclidean plane. An origami
realization of \(T_\Lambda\) is an embedding in \(\R^3\) that is
affine and distance preserving on each triangle of a triangulation of
\(T_\Lambda\). Intrinsic distances here are Euclidean distances in a
planar lift of the triangle. Equivalently, the embedded surface
inherits its metric by gluing the Euclidean triangles of the spatial
realization. It is flat precisely when all vertex angle sums are
\(2\pi\). We use the term \emph{universal} for a fixed triangulation
that realizes \(T_\Lambda\) for every rank-two lattice \(\Lambda\).

Burago and Zalgaller's isometric embedding theorem for
polyhedral surfaces already implies that every flat torus has an
origami realization after subdivision \cite{bz}. Tsuboi gave an
explicit construction for every modulus using the geometry of
diplotori \cite{tsuboi}. Lazarus and Tallerie constructed a single
universal triangulation with 2,434 faces \cite{lt}. Quintanar gave
exact coordinates for a square torus with 40 vertices
\cite{quintanar}. These results predate the classification at the
minimum vertex count.

Schwartz established the minimum vertex count of eight
\cite{schwartz}. The census of small triangulations by Lutz and
Sulanke gives seven types at this count \cite{sl}. Hougardy, Lutz,
and Zelke gave integer-coordinate spatial embeddings of all
seven with unrestricted cone angles \cite{hlz}. Lander observed that the
degree-three vertex in Type \(1\) can be erased in any flat embedding,
which would leave a forbidden seven-vertex origami torus. He found
numerical examples on the other types except Type \(2\), and asked
whether Type \(2\) admits any origami realization \cite[\S6]{lander}.
Doyle and Schwartz proved near universality for Type \(7\), with
the rectangular and sharp-rhombic rays omitted by their construction
\cite{ds}. Lander then established a Type \(7\) realization of the
square torus and local coverage around it, as well as a hexagonal
realization on Type \(3\) \cite{lander}. Doyle, Lander, Schwartz,
and Trettel proved that Type \(3\) is universal \cite[Theorem~1.1]{dlst}.

Write \(K_j\) for the type with index \(j\), with the face lists in
Appendix~\ref{sec:facelists}. If \(p=(p_i)\) is a spatial placement,
\(|K_j|\to\R^3\) denotes its affine extension in barycentric
coordinates. An embedded placement means that this extension is
injective. Compactness then makes it a closed embedding. To state the
positive results, normalize the period basis to \(1,\tau=x+iy\) up to
a lattice similarity and possibly reflection, with
\[
 D=\{x+iy: 0\le x\le \tfrac12,\ y>0,\ x^2+y^2\ge1\}.
\]
One can choose a shortest nonzero period as the first generator and
reduce the real component of the second generator modulo the first.
Set \(D_{<1/2}=\{\tau\in D:x<1/2\}\). This set contains the
rectangular ray and the unit-circle arc with the hexagonal endpoint
excluded.

\begin{theorem}\label{thm:type1}
Every embedded spatial placement of \(K_1\) has a vertex whose cone
angle differs from \(2\pi\).
\end{theorem}

\begin{theorem}\label{thm:type2}
Every embedded spatial placement of \(K_2\) has a vertex whose cone
angle differs from \(2\pi\).
\end{theorem}

\begin{theorem}[Doyle, Lander, Schwartz, and Trettel]\label{thm:type3}
The triangulation \(K_3\) realizes \(T_\Lambda\) for every rank-two
lattice \(\Lambda\).
\end{theorem}

\begin{theorem}\label{thm:type6}
The triangulation \(K_6\) realizes \(T_\Lambda\) for every rank-two
lattice \(\Lambda\). The embedding can be chosen so that the spatial
map to seven cone angles and the two real components of a locally
developed marked modulus has rank nine at its vertex coordinates.
\end{theorem}

\begin{theorem}\label{thm:type4}
For every \(\tau\in D_{<1/2}\), the triangulation \(K_4\) realizes
\(T_{\Z+\Z\tau}\), with rank nine for the spatial angle and marked
modulus map.
\end{theorem}

\begin{theorem}\label{thm:type7}
For every \(\tau\in D_{<1/2}\), the triangulation \(K_7\) realizes
\(T_{\Z+\Z\tau}\), with rank nine for the spatial angle and marked
modulus map.
\end{theorem}

In particular, \(K_7\) realizes every rectangular lattice torus,
with arbitrary scale.

\begin{theorem}\label{thm:type5square}
The triangulation \(K_5\) admits a realization of \(T_{\Z+\Z i}\) whose
coordinates are given by rational arithmetic and square roots.
\end{theorem}

\begin{theorem}\label{thm:type5box}
Let \(x_0=0.174241221618834\) and \(y_0=1.084897889755590\), interpreted
as exact rationals. The triangulation \(K_5\) realizes \(T_{\Z+\Z(x+iy)}\)
for every \((x,y)\) in the closed box
\[
 |x-x_0|\le 10^{-11},\qquad |y-y_0|\le 10^{-11}.
\]
\end{theorem}

Together the theorems settle existence for all seven types. By
Theorem~\ref{thm:type2}, the existence conjecture raised for \(K_2\)
has a negative answer. Our proof of Theorem~\ref{thm:type6} identifies
the type by its full face list. The preprint \cite[\S6]{dlst} reports
a second exactly checked universal construction by degree sequence
\(55666677\), shared by \(K_4\) and \(K_6\). That description leaves
its relation to the construction here unresolved.

\section{Nonexistence from determinant signs}\label{sec:obstruction}

\subsection*{Hemisphere support for Type 2}

Use the relabeled face list
\[
 046,\ 047,\ 056,\ 057,\ 135,\ 136,\ 145,\ 147,
 \ 167,\ 234,\ 237,\ 245,\ 256,\ 267,\ 346,\ 357.
\]
The map from the labels in Appendix~\ref{sec:facelists} is
\((0,4,7,6,5,1,3,2)\). The graph is \(K_8\) with \(01,02,03,12\)
removed. In particular \(4,5,6,7\) are universally adjacent.

Suppose \(p\) is an embedded flat placement. Intersect the star at
\(v\) with a small sphere about \(p_v\) and rescale to the unit sphere.
The link is a simple polygon of shortest great-circle arcs. Its
vertices are the unit edge directions from \(v\), and its length is
\(2\pi\). We claim that the directions are contained in an open
hemisphere. If their convex hull contains the origin, take a subset
minimal for this property. It has at most four vectors by
Carath\'eodory's theorem. The convex relation has positive coefficients,
and every proper subset is linearly independent. Otherwise one could
subtract a multiple of a proper dependence until a coefficient first
vanished and then renormalize.

For four vectors \(a,b,c,d\) in link order the positive relation gives
antipodes \(r,-r\) in the interiors of arcs \(ab,cd\), respectively.
The perimeter of the ordered quadrilateral is the sum of the lengths
of the broken paths \(rbc(-r)\) and \((-r)dar\), each at least \(\pi\).
Equality for the first forces \(r,b,c\) into one plane through the
origin and hence forces \(a,b,c\) into that plane. The perimeter
therefore exceeds \(2\pi\). Shortening the link subpaths between the
four chosen directions gives the opposite inequality.

For three vectors their short arcs partition a great circle. Each
link subpath between successive chosen vectors must have minimal
length, so the star is planar. Its vertex set contains a complete graph
on five vertices. For \(v\in\{0,1,2,3\}\) take \(v,4,5,6,7\), and
otherwise take \(3,4,5,6,7\). A planar straight-line embedding of
\(K_5\) is impossible. Three collinear vertices already force an
edge overlap. With no such triple, a hull of at least four vertices
has crossing diagonals, while a triangular hull has two interior
vertices. The three segments from the first interior vertex to the
corners partition the hull into three triangles. The segment from
the second interior vertex to the corner opposite its triangle
crosses an edge of the partition.

For two vectors the directions of neighbors \(a,b\) are antipodal.
The two paths between them on the link have length at least \(\pi\)
each, and hence both are minimizing semicircles. An edge \(ab\)
would contain \(p_v\), so \(a,b\in\{0,1,2,3\}\). At least two of the
universally adjacent neighbors have rays in the interior of one
semicircle. Write their order as \(a,w,x,b\). The ray through \(w\)
meets the relative interior of segment \(ax\), and \(vw\) must end
before that intersection. Thus \(p_w\) has strictly smaller height
above the line \(p_ap_b\) in this half-plane than \(p_x\). Applying
the same argument to \(vx\) and \(wb\) reverses the inequality.
All minimal sets have now been excluded. Strict separation of the
origin from the convex hull yields a supporting normal with positive
dot product against every neighbor direction.

\subsection*{A face with neighbors on both sides}

Radially project a supported link to the plane \(n\cdot q=1\), where
\(n\) is the supporting normal. The arcs project to segments and the
link stays simple. If every incident face plane has the other
neighbors on one closed side, every projected link edge lies on the
boundary of the convex hull. The link is a simple closed curve, so
it must be the whole hull boundary. A proper subset would lie in a
punctured convex polygon boundary, which is homeomorphic to an open
interval.

Enclose the planar hull in a triangle and pull it back to the
sphere. If \(a,b,c\) are the spherical corners, the perimeter is
\[
 2\pi + d(b,c)-d(-a,b)-d(-a,c)<2\pi.
\]
The strict inequality follows because \(-a\) lies outside the open
hemisphere, whereas the unique shorter arc from \(b\) to \(c\) lies
inside. Clip the enclosing triangle by the supporting half-planes
of the hull. Each cut removes a boundary path and replaces it, on
the sphere, by a shortest arc, which cannot increase angular
perimeter. The final perimeter is therefore less than \(2\pi\),
contrary to flatness. For the sign system we use one such mixed-side
face at each of \(0,1,2\).

\subsection*{From geometry to a finite contradiction}

Support, the selected strict side inequalities, and embeddedness
persist under sufficiently small perturbation. For embeddedness,
take two barycentric weight vectors on faces with disjoint supports.
The pairs of such vectors form a compact set defined by closed simplex
constraints and coordinatewise zero products. The minimum separation
of their images is positive, and there are finitely many face pairs.
Subtracting common barycentric weights and renormalizing reduces any
unintended coincidence to one with disjoint supports. Moving
vertex \(i\) by \(t(i,i^2,i^3)\) makes each affine four-point
determinant a polynomial in \(t\) with nonzero cubic coefficient by
the Vandermonde formula. We can avoid the finitely many roots while
retaining the strict conditions. We then reflect if necessary to
make \([0123]>0\), where \([abcd]\) denotes the determinant with
rows \((1,p_i)\) in the listed order.

The Boolean variables are the signs of the 70 sorted determinants
and 14 selectors for the mixed-side faces. Alternation supplies the
sign of every reordering. Schewe's SAT formulation of ordinary
polyhedral realizability is a precedent for the chirotope and
edge--face clauses \cite[\S2.4]{schewe}. The flatness argument here
supplies the strict support and mixed-side constraints. The clause
generator uses the following conditions.

\begin{enumerate}
\item For \(S\) an increasing pair disjoint from \(a<b<c<d\), the
three signed terms of
\[
 [Sab][Scd]-[Sac][Sbd]+[Sad][Sbc]=0
\]
cannot all have the same sign. This is the two-dimensional
Grassmann--Pl\"ucker identity after quotienting by the rows of \(S\).
\item For an increasing five-tuple \(C\), the numbers
\((-1)^i[C\setminus C_i]\) are the coefficients of an affine
dependence, with \(i\) starting at zero. They have both signs.
After normalization the positive and negative parts are equal
convex combinations. A sign partition into an edge and a disjoint
face contradicts embeddedness. A singleton and four of its
neighbors contradict strict support at the singleton.
\item At each of \(0,1,2\) at least one selector is true. Its face
\(abc\) has brackets of both signs among \([abcw]\) for neighboring
\(w\) outside the face.
\end{enumerate}

Including normalization, this gives 7,464 clauses. The refutation
uses 741 of them and 363 reverse-unit-propagation steps. For each new
clause, temporarily falsify every literal of that clause. Scan its
listed predecessor clauses in order. A single remaining literal
forces that literal to be true, and a clause with no remaining
literal produces a contradiction. Each such contradiction proves the
new clause. The final new clause is empty. As an example, the first
derived clause is \((-61,18,-39,-19,-38)\) in signed-literal
notation. The input unit clause \((1)\) and the input clause formed
by adjoining \(-1\) to the new clause give the replay in this case.

The replay and the indexing of its inputs appear in the
\src{OrigamiTorus/Type2/Certificate.lean}{Type 2 certificate}.
Lean checks the connection to the geometric clauses and the replay
itself. The perturbation is only a means of obtaining strict signs
and imposes no genericity requirement on the original flat
placement. This proves Theorem~\ref{thm:type2}.

\subsection*{Type 1 and the collapsed seven-vertex complex}

At the degree-three vertex, opposite directions would put the edge
between their endpoints through the center of the star. At flatness
the three link arcs therefore have length sum \(2\pi\) with each
length strictly less than \(\pi\). The third direction lies on a
shortest path between the antipodes of the first two, by equality
in the angular triangle inequality. The three directions thus lie
on a great circle and positively span its plane. The star center
lies strictly inside the triangle of its neighbors. Replacing the
three incident triangles by that triangle preserves the image and
all the other cone angles.

The graph on the seven remaining vertices is complete. The same
minimal-set proof gives hemisphere support at every remaining
vertex. The two-vector case is excluded directly by completeness.
The other cases use the planar \(K_5\) obstruction and the
four-direction perimeter inequality. Generic perturbation yields
35 sorted determinant signs satisfying the circuit and
Grassmann--Pl\"ucker restrictions. On the particular collapsed
face list these give 2,100 clauses. Lean verifies the 733 input
clauses used in a 418-step refutation in
\src{OrigamiTorus/Type1/Certificate.lean}{Type1/Certificate.lean}.
This gives a proof of Theorem~\ref{thm:type1} on the collapsed complex
and formalizes the obstruction needed by Lander's star reduction.
Schwartz's theorem concerns all seven-vertex paper tori
\cite{schwartz}.

\section{From planar data to a spatial realization}\label{sec:method}

Two tests recur in the constructions. The first shows that a strict
height lift embeds the spatial face complex, and the second
identifies the developed intrinsic faces with the requested flat
torus.

\subsection*{Separating a disjoint edge and face}

For an ordered edge \(ab\) and a disjoint face \(cde\), let \(V_i\)
be the determinant of the ordered list \((a,b,c,d,e)\) with position
\(i\) deleted. In general position the affine dependence on these
five points has coefficients proportional to
\((V_0,-V_1,V_2,-V_3,V_4)\). If the segment and triangle meet, the
coefficients of the edge must have one sign and those of the triangle
the other. Thus any of
\[
 V_0V_1>0,\qquad V_2V_3>0,\qquad V_3V_4>0,\qquad V_2V_4<0
 \tag{1}\label{eq:separation}
\]
excludes intersection. For two arbitrary faces, subtract common
barycentric mass from a hypothetical unwanted intersection, leaving
two disjoint supports of size at most three. With six nonzero weights
there is a nonzero variation of the weights preserving the two
weight sums and the equality in \(\R^3\). Move until one weight first
vanishes. General position then forces exactly five points in the
dependence, with support sizes two and three. It suffices to check
\eqref{eq:separation} on all disjoint edge--face pairs. If the
strict tests hold at some height vector and every projected
four-tuple contains a noncollinear triple, a small change of heights
can avoid all zero affine four-point determinants. The strict tests
survive this change.

\subsection*{Positive periodic triangles and area}

Let a collection of oriented planar lifts of the abstract faces
have compatible edge pairings by lattice translates. Suppose all
oriented face areas are positive and sum to \(\covol(\Lambda)\).
The resulting map \(|K|\to T_\Lambda\) is open. In the interior of
a face it is locally affine invertible. At an edge the two positive
triangles fill opposite sides locally. At a vertex all incident
triangles can be translated to a common central corner. Their
consecutive oriented sectors sweep positive angles less than \(\pi\),
and because the cycle closes these angles sum to a positive multiple
of \(2\pi\). Every direction therefore occurs in a sector. The
finite fan maps every neighborhood of the vertex onto a
neighborhood of its image. The image is compact, nonempty and
open in the connected torus, hence surjective.

Cut the translates of the lifted triangles along a fundamental
parallelogram. This gives finitely many pieces whose total area
equals the area of the cell. If a point had two distinct preimages,
openness would give an open set of overlap between the images of
disjoint neighborhoods of the preimages. A point there off all
projected edges and off the cell boundary would lie in the interiors
of two translated triangles. Their positive-area overlap contradicts
equality of the sum of the piece areas with the area of their
union. Thus the map is a homeomorphism. Equality of the three edge
lengths of each spatial face with its planar lift makes the affine
map on that face an isometry. This argument gives a realization with
the prescribed lattice, at the prescribed scale.

\section{The universal constructions}\label{sec:universality}

\subsection*{Planar exchange}

For the two charts of \cite[\S3]{dlst}, set \(w=y^2\) and choose
\[
 (u,k,\lambda)=
 \begin{cases}
 (3/128,11/32,13/32),&0\le x\le1/6,\\
 (x/2-1/2048,1/16,1/2),&1/6\le x\le1/2.
 \end{cases}
\]
At \(x=1/6\) either chart applies. Put \(b=1-x-2k\).
The folded planar points have coordinates \((X,Y/y)\) with the
following scaled coordinates \((X,Y)\).
\begin{center}
\begin{tabular}{@{}cll@{}}
\toprule
\(i\)&\(X_i\)&\(Y_i\)\\
\midrule
0 & \(0\) & \(0\)\\
1 & \(u\) & \(u(b+2k/\lambda)\)\\
2 & \(u-1/2\) & \((1-2u)x/2\)\\
3 & \(-k-b\lambda\) & \((\lambda(w-b^2)-2bk)/2\)\\
4 & \(x+k-1/2\) & \((w+bx)/2\)\\
5 & \(u-1/2\) & \((w+b(b+2u-1))/2\)\\
6 & \(u-k\) & \(b(u-k)\)\\
7 & \(u\) & \((w+b^2+2bu)/2-(2u-2x-2k+1)/(2(1-\lambda))\)\\
\bottomrule
\end{tabular}
\end{center}
The planar intrinsic representatives \(Q_i\) agree with the
folded representatives except at \(3\) and \(6\), where their
scaled coordinates are
\[
 Q_3=\bigl(k,-(2bk+\lambda(b^2+w))/2\bigr),\qquad
 Q_6=\bigl(u+k,b(u-k)\bigr).
\]
The directed period shifts are listed in
Appendix~\ref{sec:shifts}.

In both charts \(0<u<1/2\) and \(b>0\). The points \(2,5\) have the
same first coordinate, and \(Y_5>Y_2\), since
\[
 2(Y_5-Y_2)=w+b(b-1)-x+2u(b+x)>0.
\]
Indeed, \(w+b(b-1)-x\ge 3/4-x^2-x\ge0\) and \(b+x=1-2k>0\).
Also \(Y_1>0\) and \(Y_5>Y_2\ge0\). Taking determinants therefore
gives the signs \(+,-,-,+\) for \(120,125,051,052\).
The diagonals \(12\) and \(05\) cross in the interior of
the convex quadrilateral \(0,1,5,2\). The original faces \(012\)
and \(152\) lift with zero shifts. Replacing them with \(015\) and
\(052\) retriangulates the same quadrilateral on the intrinsic
sheet and on the folded sheet. The new edge \(05\) was absent in the
original face list. Relabeling by \((3,7,0,6,1,5,2,4)\) gives \(K_6\).

\subsection*{Height and length certificates}

The Type \(6\) height vectors \(z\), in chart order, are
\[
 \begin{aligned}
 &(0,0,-11,-1,-3,-15,0,977),\\
 &(0,-361,2527,2525,5105,2527,1,19639).
 \end{aligned}
\]
The complementary minor signs give at least one strict test in
\eqref{eq:separation} for each of the 140 disjoint edge--face pairs
in each chart. To check the arithmetic, write
\(x=l(1-s)+rs\) on a chart interval \([l,r]\) and
\(w=1-x^2+h\) with \(h\ge0\). If \(A_{abc}\) is the oriented
double area of a projected triple, a volume is
\[
 -z_a A_{bcd}+z_b A_{acd}-z_c A_{abd}+z_d A_{abc}.
\]
Write the signed area or volume as a cubic in \(1-s,s\), plus
\(h\) times a linear polynomial in \(1-s,s\). Positive coefficients
of the cubic and nonnegative coefficients of the linear polynomial
prove the claimed strict sign on the closed strip. The sign table,
including permutation parity, is in
\src{OrigamiTorus/Type6/Signs.lean}{Type6/Signs.lean}.
The 70 volumes for the first chart's vector are nonzero.
In the second chart, the auxiliary vector
\((0,-2017,0,-2,-1,-3,0,1430)\) has all 70 volumes nonzero.
Adding a sufficiently small multiple of that vector to the second
chart's lift preserves strict separation and avoids coplanarity.

Hold the periods fixed. Write \(P_i\) and \(Q_i\) for the movable
folded and intrinsic horizontal coordinates in the scaled
plane. For a directed edge \(ij\) with shift \((m,n)\), put
\(d_{ij}=(m+nx,nw)\) in scaled coordinates and
\(\|v\|_w^2=w v_X^2+v_Y^2\). The residual for that edge is
\[
 \|P_j-P_i\|_w^2+w t^2(z_j-z_i)^2
       -\|Q_j+d_{ij}-Q_i\|_w^2.
 \tag{2}\label{eq:residual}
\]
There are 24 residuals. They vanish at the planar base and \(t=0\).
Choose both coordinates of folded vertices \(0,\ldots,5\) and the
first at \(6\), and both coordinates of intrinsic vertices
\(0,\ldots,3,5\) and the first at \(4\). Each derivative row at the
base is twice the difference of the weighted dot products of the
folded and intrinsic edge vectors with the endpoint motions.
The coefficient elimination in
\src{OrigamiTorus/Type6/Rank.lean}{Type6/Rank.lean}
proves that the square derivative has trivial kernel on both
charts. It uses \(h=w-1+x^2\ge0\) to justify every nonzero factor
canceled in the row operations.

The implicit function theorem gives horizontal positions near the
base at every sufficiently small \(t>0\) with all residuals zero.
After dividing the vertical coordinate by \(t\), these positions
approach the strictly embedded lift. By openness of the embedding
condition, the adjusted lift is embedded and stays embedded when
the vertical scale is restored. The intrinsic triangles remain
positively oriented. The period table gives opposite edge pairings
and total area equal to one period cell, with twice the scaled
area equal to \(2w\). The area argument of
Section~\ref{sec:method} gives the required lattice realization.
Lattice reduction and similarity transport it to any \(\Lambda\),
rescaling the spatial vertices with the same metric scale.

\subsection*{Regularity at the Type 6 embedding}

The two marked periods can be recovered locally by successive
Euclidean triangle developments along
\[
 \begin{aligned}
 &167\to276\to257\to052\to015,\\
 &167\to276\to257\to375\to356\to136.
 \end{aligned}
\]
Normalize the common starting face by placing vertex \(1\) at \(0\)
and vertex \(7\) at \(1\). The two paths end at the \((1,0)\) and
\((0,1)\) lifts of vertex \(1\), respectively. Their displacement
ratio equals the period ratio on any nearby positive periodic
sheet. Thus the marked modulus is a smooth function of the squared
edge lengths near the constructed metric.

To see the angle directions, vary squared lengths \(L_{ij}\) by
\(\dot L_{ij}=(u_i+u_j)L_{ij}\). In a nondegenerate triangle the
cosine rule gives
\[
 \dot A_i=\tfrac12\bigl((u_j-u_i)\cot A_k+
                         (u_k-u_i)\cot A_j\bigr).
\]
The cotangent quadratic form \(-\sum_i u_i\dot A_i\) is the
Euclidean vertex-scaling Hessian up to normalization
\cite[Propositions~4.1.6 and~4.2.3]{bps}. It is positive on the
quotient by constant vertex weights. This follows by writing the
form using the Gram matrix of two incident sides, whose determinant
is positive.
Summing over faces and using connectedness gives seven
independent cone-angle variations. Varying the two coordinates of
the second period of the intrinsic sheet gives two modulus
directions with all angles fixed, since the positive periodic
faces continue to tile their quotient. Thus the derivative \(D\)
of the joint angle and modulus map on squared lengths is onto.
At the adjusted positive-height solution the derivative of
\eqref{eq:residual} in the chosen coordinates is still invertible.
Write it in physical units as \(\dd E-\dd I\), where \(E,I\)
are the folded and intrinsic squared lengths. Since the intrinsic
variations keep the periods fixed, \(D\,\dd I=0\).
Surjectivity of \(D(\dd E-\dd I)\) gives surjectivity of
\(D\,\dd E\) at the final spatial vertices. Similarity preserves
this rank. This proves Theorem~\ref{thm:type6}.

\subsection*{The Type 3 check}

For the original diagonal \(12\), the formalization uses the
published Type \(3\) heights \cite[\S4]{dlst}. Here is how the two
length derivatives are related. For the convex quadrilateral,
take a nonzero affine dependence \((\alpha_i)\), with
\(\sum\alpha_i=0\) and \(\sum\alpha_iP_i=0\). No coefficient is
zero. The stress relation on its six chords is
\[
 \sum_{i<j}\alpha_i\alpha_j
       \langle P_i-P_j,\dot P_i-\dot P_j\rangle_w=0.
\]
It follows by expansion and holds on the intrinsic sheet with
the same coefficients, because these four representatives agree
there and all six relevant shifts are zero. Subtract the two
relations. The coefficient of each diagonal is nonzero, so
replacing the derivative row for \(05\) by the row for \(12\)
preserves invertibility. The same fixed-period inflation and area
argument, now using the original positively oriented face lifts,
verify Theorem~\ref{thm:type3}. The stress argument and quotient
proof are given in
\src{OrigamiTorus/Type3/Rank.lean}{Type3/Rank.lean} and
\src{OrigamiTorus/Type3/Universal.lean}{Type3/Universal.lean}.

\section{The trapezoidal constructions for Types 4 and 7}\label{sec:strip}

Use the unflipped \(K_7\) face list in Appendix~\ref{sec:facelists}.
For \(w=y^2\), the two charts use
\[
 (H,c)=
 \begin{cases}
 (33/(20+26x),1/3),&0\le x\le9/20,\\
 (1+3(1-2x)^2/2,2/3),&9/20\le x<1/2.
 \end{cases}
\]
In the planar coordinates \((X,Y)\) used below the physical point
is \((X,yY)\). Put
\[
 \begin{gathered}
 a=Hx,\quad D_0=H-1+a,\quad Q=D_0^2+H^2w,\quad
 q=H/D_0,\quad r=-cD_0,\\
 T=a^2-a+H^2w-D_0,\qquad
 (1-c)S=a^2-a-(1-c)H^2w-D_0(2a-cD_0-1),\\
 s=D_0T/Q,\qquad t=D_0S/Q.
 \end{gathered}
\]
The eight points in label order are
\[
 \begin{array}{llll}
 (0,0),&(1-a-r,H+qr),&(1/2,0),&(a+t,H+qt),\\
 (1,0),&(1+s,qs),&(a,H),&(1-a,H).
 \end{array}
\]
The bound \(a^2+H^2w\ge H^2\) gives
\(T\ge(H-1)^2+H(1-2x)>0\) and
\[
 (1-c)S\le (H-1)+c(H-1)^2-(1-c)H^2-2(1-c)(H-1)a<0.
\]
The last inequality uses \(H^2-H+2>0\) in the first chart
and \(1<H\le 6/5\) in the second. Also
\[
 (1-c)\bigl(Q+(2H-1)T+S\bigr)=(1-c)(2H-1)T+H(H-1)>0.
\]
The double areas of \(356,352,623,625\) after scale \(Q\) are
\[
 \begin{aligned}
 &-SH^2Q,\qquad &&-HQ\bigl(Q+(2H-1)T+S\bigr)/2,\\
 &SHQ(H-1/2), &&HQ\bigl(Q+(2H-1)T\bigr)/2.
 \end{aligned}
\]
Their signs are \(+,-,-,+\). Thus \(35\) can be exchanged for \(26\)
in a convex quadrilateral. Relabeling by \((2,1,4,0,3,7,5,6)\)
gives \(K_4\).

The angle and period calculations use two circles. Vertices
\(0,4,5,6,7\) are on the circle whose equation in these coordinates
is
\[
 X^2+wY^2-X-\frac{a^2-a+H^2w}{H}Y=0.
\]
Vertices \(0,4,1,3\) lie on a second circle, along with the
reflection of \(1\) in the vertical line \(X=1/2\). For the
oriented link at any vertex, start with the successive ratios of
its complex edge directions. They telescope over the full link.
In passing to positive interior angles, reverse the phase for
each negatively oriented face. The ratios over a consecutive
negative block telescope to its end chords. The two circle
identities and the crease directions make the product of the
negative-block ratios real at each vertex. The corrected phases
therefore multiply to one. Every cone angle is a positive
multiple of \(2\pi\), and the sixteen faces have total angle
\(16\pi\), so every cone angle is \(2\pi\). The flip preserves
these angle sums. The full chord calculation is in
\src{docs/Type4.md}{docs/Type4.md}.

For the periods, unfold the six faces adjacent to the marking
paths \(4\to2\to3\to4\) and \(5\to2\to0\to5\). The loops meet at
\(2\) with alternating branches, and the flip preserves both
loops. Put \(\rho=(D_0,H)\) and \(\nu=(-wH,D_0)\).
The creases encountered in the two continuations have directions
\(\rho\) and \(x\rho+\nu\). Applying reflection across these
directions in the metric \(X^2+wY^2\), the two developed vectors
in the common scale \(Q\) are
\[
 U=2H^2\nu,\qquad V=2H^2(x\nu-w\rho).
\]
The vectors \(\nu,\rho\) are orthogonal in that metric and
\(\|\nu\|^2=wQ\). The physical ratio \(V/U\) is \(x+iy\).
The signed corner derivatives have seven independent angle
directions. In their left kernel the weights at the eight
vertices satisfy linear balance equations in the inverse complex
chords at the negative-block boundaries. Projecting the balance
equations at vertices \(3,5\) along the two circles, then
eliminating using the equations at \(1,6\), gives only the
constant weights. The two chart derivatives preserve cone angles
and vary the modulus independently. The resulting horizontal rank
is nine. The flip preserves the angle map in a planar
neighborhood of the base, and both planar charts give the same
marked ratio before and after the flip. Lean checks the rank
calculation in \src{OrigamiTorus/Type4/StripRank.lean}{Type4/StripRank.lean}.

There are six height cells, cut at \(0,1/10,1/5,3/10,2/5,9/20,1/2\).
The first five are closed and the last excludes \(1/2\).
Use \(w=1-x^2+h\) and a coordinate \(s\in[0,1]\) on each cell.
The coordinate and height data give determinant polynomials
after clearing positive denominators. On \(s<1\), substituting
\(s=v/(1+v)\) and clearing powers of \(1+v\) reduces the required
strict signs to nonnegative coefficients and a strictly positive
constant coefficient for \(v,h\ge0\). On a closed cell the
endpoint \(s=1\) is checked separately. The checks cover all
disjoint edge--face pairs of each face list. For \(K_7\) the second
cell uses the first height polynomial evaluated at \(1+s\),
and the third adds \((1+h)/100\) to the cleared height of vertex
\(3\). The other cells share the \(K_4\) height rows.
The finite data are in
\src{OrigamiTorus/Type4/LiftWideData.lean}{Type4/LiftWideData.lean}
and \src{OrigamiTorus/Type7/WideData.lean}{Type7/WideData.lean}.

With these embedded lift directions, the rank-nine angular and
marked-modulus equations admit an implicit-function correction at
small positive height. The seven angle equations force the
eighth by the face-angle total. The specified marked ratio is
fixed during the correction. The full spatial derivative already
has rank nine at the planar base, so this rank persists at
sufficiently small positive height and at the final rescaling.

To identify the corrected metric, successive oriented unfoldings
give face transitions in the orientation-preserving Euclidean
group. The angle equations make the transitions around each
vertex trivial. After gauging a spanning tree in the dual
graph, the other transitions are words in two commuting
generators. For example, on the unflipped list write \(a,b,c,d,e,f,g,u,v\)
for the transitions across \(01,02,03,05,06,07,14,12,47\)
in that gauge. The eight vertex words, in label order, are
\[
 \begin{gathered}
 bcaedf,\quad gua^{-1},\quad b^{-1}u^{-1},\quad c^{-1},\\
 vg^{-1},\quad d^{-1},\quad e^{-1},\quad f^{-1}v^{-1}.
 \end{gathered}
\]
Setting them to the identity gives \(c=d=e=1\) and
\((a,b,f,g)=(vu,u^{-1},v^{-1},v)\), with \(uv=vu\).
If either generator is a nontrivial rotation, every commuting
transition fixes its center. Summing the boundary determinant
formula for the oriented areas about that center cancels every
paired edge, contradicting positive total area. Thus all
transitions are translations. In the unflipped gauge the marked
periods read off the translations of \(v\) and \(uv\), a unimodular
change of basis. The shifted face lifts have total area equal to
one period cell. The checked developments and area identities for
both lists are in
\src{OrigamiTorus/Type7/IntrinsicLifts.lean}{Type7/IntrinsicLifts.lean}
and \src{OrigamiTorus/Type7/Tiling.lean}{Type7/Tiling.lean}.
The quotient argument finishes the realizations in
Theorems~\ref{thm:type4} and~\ref{thm:type7}.

\section{Exact radical realizations of Type 5}\label{sec:radicals}

Work in the labeling of Appendix~\ref{sec:facelists}. The anchors
are \(1,2,5\), the roots are \(3,4\), and the leaves are \(0,6,7\).
There is no edge between the two roots, and no edge joins two
leaves. Every leaf is adjacent to all five core vertices.
Use the integral shifts \(K_{ij}\) given in
Appendix~\ref{sec:shifts}, interpreted in the requested basis.
Vertex \(1\) is the origin and all shifts incident to it vanish.
The anchors lie in the horizontal plane at their respective
intrinsic positions. The following integers, divided by
\(10^{15}\), give the physical planar core coordinates.
\begin{center}
\small
\begin{tabular}{@{}crrrr@{}}
\toprule
&\multicolumn{2}{c}{Square}&\multicolumn{2}{c}{Nonrectangular box}\\
\(i\)&\(10^{15}X_i\)&\(10^{15}Y_i\)&\(10^{15}X_i\)&\(10^{15}Y_i\)\\
\midrule
2 & \(444720882667740\) & \(-499972918103549\) & \(305146237868087\) & \(-605447047939750\)\\
3 & \(499999999999045\) & \(87262616850970\) & \(499997454165673\) & \(181480875778771\)\\
4 & \(-81664706988789\) & \(418104802964447\) & \(-100309264271995\) & \(459493147984820\)\\
5 & \(-39101278348661\) & \(-213154941393696\) & \(-154596026157111\) & \(-274376548947718\)\\
\bottomrule
\end{tabular}
\end{center}

To reconstruct a vertex \(j\), write \(u\in\R^2\) for its
intrinsic position, \(v\in\R^3\) for its spatial position, and
\(r_i\) for the spatial position of a core neighbor \(i\).
Set \(s_i=p_i-K_{ij}\), with the shift understood as a physical
plane vector. Subtract the equation \(\|v\|^2=\|u\|^2\) from the
squared-distance equation for each other neighbor to obtain
\[
 r_i\cdot v-s_i\cdot u=(\|r_i\|^2-\|s_i\|^2)/2.
 \tag{3}\label{eq:radical}
\]
The equations for anchors \(2,5\) determine the horizontal part
of \(v\) by a \(2\times2\) solve. At a root, \(u\) is given, and
the height is a signed square root of the remaining squared norm.
Take the positive sign at \(3\) and the negative sign at \(4\).

At a leaf the horizontal part is \(B+Lu\) for a vector \(B\)
and matrix \(L\). Substitution in the two root equations gives
\(z_i v_z=a_i\cdot u+b_i\) for \(i=3,4\). Eliminate \(v_z\),
giving \(n\cdot u+m=0\), where
\[
 n=z_4a_3-z_3a_4,\qquad m=z_4b_3-z_3b_4 .
\]
The checked data have \(n\ne0\) and \(z_3\ne0\). Put
\[
 u(t)=-\frac{m}{\|n\|^2}n+t(-n_y,n_x),\qquad
 v(t)=\left(B+Lu(t),\,\frac{a_3\cdot u(t)+b_3}{z_3}\right).
\]
Solving the quadratic \(\|v(t)\|^2-\|u(t)\|^2=0\) supplies the
leaf coordinates. Use the plus sign before the square root
in the quadratic formula for \(0,7\), and the minus sign for
\(6\). Each reconstructed leaf satisfies all five incident
lengths. Together with the two roots and the anchor triangle,
this accounts for all 24 edges.

For the square and the entire nonrectangular box, Lean verifies
the necessary nonzero denominators, positive radicands, 16
positive face orientations, 70 nonzero four-point determinants,
and strict tests on the 138 disjoint edge--face pairs.
Dyadic and rational interval bounds enclose the radicals and
prove these strict comparisons. Exact identities, including
\eqref{eq:radical}, prove the edge equalities.
The skew shifts have zero sum around each face and the lifts
have compatible edge pairings. The sum of signed double areas
is 2 in lattice coordinates, since all terms depending on the
vertex representatives cancel in pairs. Physical coordinates
multiply this area by \(y\). The argument in
Section~\ref{sec:method} gives the desired square and nonrectangular
quotients. The formulas and interval checks are in
\src{OrigamiTorus/Type5/Square.lean}{Type5/Square.lean} and
\src{OrigamiTorus/Type5/ContinuationCertificate.lean}{Type5/ContinuationCertificate.lean}.
This proves Theorems~\ref{thm:type5square} and~\ref{thm:type5box}.

The same length elimination bounds the number of reconstruction
candidates for a fixed marking and four intrinsic core points.
For a flat embedded placement supplied with coherent,
nondegenerate marked face lifts, the anchor triangle can be
put in this gauge by a rigid motion. The formal proof excludes
a rank loss in the leaf equations and an identically zero
remaining norm polynomial. There are then at most two choices
at each root and at each leaf, giving at most \(2^5=32\)
candidates. Some candidates can fail the embedding and face
conditions. The completeness statement takes the marked lifts
as input. See
\src{OrigamiTorus/Type5/MarkedReduction.lean}{Type5/MarkedReduction.lean}.

\section{Remaining universality questions}

Within the reduced domain, the remaining moduli for \(K_4\) and
\(K_7\) lie on the sharp-rhombic ray \(x=1/2\),
\(y\ge\sqrt3/2\). The covered regions for \(K_5\) are smaller.
One can restrict the possible markings of a realization before
solving for its spatial coordinates.

For a piecewise-isometric embedded realization, an edge is
strictly shorter than every alternative period lift of the
segment with the same endpoints in the quotient. In fact,
subdivide the alternative planar segment along the translated
face lifts. Its spatial image has length at most that of the
segment and at least the Euclidean distance between the
edge endpoints. Equality forces the whole image onto that
spatial edge. Injectivity of the realization and uniqueness
of a lift with fixed starting point then force the same
lift. A finite-subdivision version also works when the
segment lies partly on lifted face boundaries.

The resulting strict nearest-period inequalities bound the
period of an essential triangular cycle to a finite set.
Using a pair of generating cycles, choose a unimodular
period matrix from these possibilities. Gauge the shifts
on a spanning tree to zero. The two cycle periods, followed
by the face equations, then determine every edge shift for
the three types. Enumerating under combinatorial automorphisms and
specified lattice isometries gives the following numbers of orbits
of necessary marking candidates.
\begin{center}
\begin{tabular}{@{}lrrr@{}}
\toprule
Type & \(x<1/2\) & Sharp ray & Hexagonal endpoint\\
\midrule
\(4\)&6&2&1\\
\(5\)&2&1&1\\
\(7\)&10&4&2\\
\bottomrule
\end{tabular}
\end{center}
The column \(x<1/2\) uses central inversion uniformly.
Additional symmetries at special moduli can identify further
candidates there. The sharp-ray column uses the symmetries of the
nonhexagonal part of the ray.
At the hexagonal metric all the ray representatives admit
positively oriented intrinsic triangulations with strictly
shortest marked edges. The intrinsic construction uses
rational coordinates in a lattice basis and the same area
test as above. This leaves a spatial embedding problem for
the four hexagonal orbits. The winding, orbit, and intrinsic
statements are in
\src{OrigamiTorus/Winding/ShortestWinding.lean}{Winding/ShortestWinding.lean},
\src{OrigamiTorus/Winding/WindingOrbits.lean}{Winding/WindingOrbits.lean},
and \src{OrigamiTorus/Winding/HexagonalTiling.lean}{Winding/HexagonalTiling.lean}.

For each fixed marking of \(K_4\) or \(K_7\), the choice of
fourteen real coordinates for the intrinsic vertex
representatives after translation fixes the 24 marked edge
lengths. The formalization lists five possible four-vertex
anchors and proves that one is a nondegenerate tetrahedron in
any embedded placement. After fixing the anchor up to rigid
motion and reflection, each remaining vertex lies on three
spheres centered at distinct anchor vertices. Subtracting the
squared-distance equations fixes its two coordinates in the
plane of the centers. The remaining coordinate is a signed
square root, including the possibility of tangency. There are at
most \(5\cdot 2^4=80\) candidates. Exact tests of all edge
lengths and disjoint subface intersections apply even when
some non-anchor vertices are coplanar. The completeness statement
takes the marking and representatives as input.
\src{OrigamiTorus/Geometry/TetrahedralReconstruction.lean}{TetrahedralReconstruction.lean}

\section{Formal verification and availability}

The Lean~4 development is available in the
\href{https://github.com/vltanh/lean4-origami-torus}{origami-torus repository}.
The results in this paper refer to the pinned
\href{https://github.com/vltanh/lean4-origami-torus/commit/\release}{source revision}.
It includes the individual proof notes, exact coordinate
tables, and complete finite certificates. The toolchain
is Lean~4.36.0 with pinned Mathlib dependencies.
The build and comparison commands run from the repository root.
{\small
\begin{verbatim}
lake build
lake env lake comparator --config=palomar/OrigamiTorus/comparator.json
\end{verbatim}
}
Lean and NanoDa accepted the 13 statements in the
Comparator configuration in a local replay. This configuration
includes the main obstructions, universality results with
Type \(6\) regularity, Type \(4\) strip existence, Type \(5\)
square, and Type \(7\) rectangular existence. The full
Lean library also checks the Type \(5\) box and the rank-nine
strip statements for Types \(4\) and \(7\). The main declarations
are as follows.
\begin{center}
\small
\begin{tabularx}{\linewidth}{@{}lX@{}}
\toprule
Result & Declaration under \texttt{OrigamiTorus}\\
\midrule
Theorem~\ref{thm:type1}&\texttt{Type1.nonexistence}\\
Theorem~\ref{thm:type2}&\texttt{Type2.type2\_nonexistence}\\
Theorem~\ref{thm:type3}&\texttt{Type3.universal}\\
Theorem~\ref{thm:type6}&\texttt{type6\_regular\_universal}\\
Theorem~\ref{thm:type4}&\texttt{Type4.open\_strip\_regular\_realizes}\\
Theorem~\ref{thm:type7}&\texttt{Type7.open\_strip\_regular\_realizes}\\
Theorem~\ref{thm:type5square}&\texttt{Type5.square\_torus}\\
Theorem~\ref{thm:type5box}&\texttt{Type5.Continuation.torus\_by\_faces}\\
\bottomrule
\end{tabularx}
\end{center}
For the square, \texttt{Type5.square\_torus} verifies the embedding
built from the radical coordinates defined in \texttt{Type5/Square.lean}.
The local comparison checks the independently stated Challenge claims
against the Solution proofs. The finite proofs replay through Lean's
kernel and require no SAT solver at proof-checking time. The project
\src{docs/VERIFICATION.md}{verification notes} record the precise
scope of the comparisons.

\subsection*{Acknowledgments}
The planar formulas underlying the Type \(3\) and Type \(6\)
constructions are due to Doyle, Lander, Schwartz, and Trettel.
The trapezoidal construction uses the geometry of the work on
near universality and the square torus. CaDiCaL assisted in
finding the Type \(2\) propositional refutation, and SymPy
assisted in finding the Type \(6\) elimination order. OpenAI
Codex (GPT-6) contributed to mathematical proof development and
the Lean implementation, carried out separate proof review runs,
and drafted the manuscript under the author's direction. The
code and proof credit are documented in
\src{docs/CREDITS.md}{docs/CREDITS.md}.

\appendix

\section{Face lists}\label{sec:facelists}

The strings \(abc\) denote triangular faces with vertex labels
\(a,b,c\). The lists below specify the triangulations regardless
of the starting point on an oriented face. They use the labeling
in the public numerical data of Lander and Trettel
\cite{data}. The proof of Theorem~\ref{thm:type2} uses the
relabeling given in Section~\ref{sec:obstruction}.

\begin{center}
\small
\begin{tabular}{@{}cl@{}}
\toprule
Type & Faces\\
\midrule
1 & \texttt{012 031 023 142 135 164 157 176}\\
  & \texttt{263 247 256 275 345 374 367 465}\\[2pt]
2 & \texttt{012 031 024 043 152 136 145 174}\\
  & \texttt{167 253 237 264 276 347 356 465}\\[2pt]
3 & \texttt{012 031 024 043 152 136 145 174}\\
  & \texttt{167 264 257 276 347 356 375 465}\\[2pt]
4 & \texttt{012 031 024 053 045 152 134 146}\\
  & \texttt{165 264 257 276 374 356 367 475}\\[2pt]
5 & \texttt{012 031 024 053 045 152 136 164}\\
  & \texttt{147 175 263 237 274 256 357 465}\\[2pt]
6 & \texttt{012 031 024 053 045 152 136 164}\\
  & \texttt{147 175 274 256 267 357 376 465}\\[2pt]
7 & \texttt{356 325 364 302 341 310 506 524}\\
  & \texttt{547 570 674 601 617 214 207 271}\\
\bottomrule
\end{tabular}
\end{center}

\section{Edge shifts for the fixed-period and radical constructions}\label{sec:shifts}

For an increasing edge \(ij\), a shift \((m,n)\) means that the
intrinsic difference is \(Q_j-Q_i+m+n\tau\) in physical complex
coordinates. In the two lists below omitted shifts are zero,
and reversing an edge negates its shift.

For the Type \(3\) and Type \(6\) charts in
Section~\ref{sec:universality} use these shifts in the original
Type \(3\) labeling, with shift zero on the relevant diagonal.
\[
\begin{array}{c|l}
(m,n)&\text{Edges}\\ \hline
(0,-1)&04,\ 24,\ 34,\ 37\\
(-1,0)&26,\ 27,\ 57\\
(1,-1)&35\\
(-1,1)&46,\ 56
\end{array}
\qquad
\begin{array}{c|l}
(m,n)&\text{Type \(5\) edges}\\ \hline
(1,-1)&04,\ 24,\ 27\\
(1,0)&05,\ 35,\ 37\\
(0,-1)&23,\ 26,\ 56\\
(0,1)&45
\end{array}
\]
For Type \(5\), the right table gives the shifts in the Type \(5\)
labeling and applies to both radical constructions. The
shift data are defined in
\src{OrigamiTorus/Defs.lean}{Defs.lean} and
\src{OrigamiTorus/Type5/Square.lean}{Type5/Square.lean}.

\begin{thebibliography}{99}

\bibitem{bps}
A. I. Bobenko, U. Pinkall, and B. A. Springborn,
\emph{Discrete conformal maps and ideal hyperbolic polyhedra},
Geometry \& Topology 19 (2015), 2155--2215.
\href{https://arxiv.org/abs/1005.2698}{arXiv:1005.2698}.

\bibitem{bz}
Yu. D. Burago and V. A. Zalgaller,
\emph{Isometric piecewise-linear embeddings of two-dimensional
manifolds with a polyhedral metric into \(\R^3\)},
Algebra i Analiz 7 (1995), 76--95.
\href{https://www.mathnet.ru/eng/aa555}{English translation}, St. Petersburg Math. J. 7
(1996), 369--385.

\bibitem{dlst}
P. Doyle, F. Lander, R. E. Schwartz, and S. Trettel,
\emph{A universal \(8\)-vertex origami torus triangulation},
preprint, 2026.
\href{https://arxiv.org/abs/2611.21909v1}{arXiv:2611.21909v1}.

\bibitem{ds}
P. Doyle and R. E. Schwartz,
\emph{Collapsibility and near universality for vertex minimal paper tori},
preprint, 2025.
\href{https://arxiv.org/abs/2510.12623v4}{arXiv:2510.12623v4}.

\bibitem{hlz}
S. Hougardy, F. H. Lutz, and M. Zelke,
\emph{Polyhedral tori with minimal coordinates},
preprint, 2007.
\href{https://arxiv.org/abs/0709.2794}{arXiv:0709.2794}.

\bibitem{lander}
F. Lander,
\emph{The two most symmetric flat tori as eight-vertex paper tori},
preprint, 2026.
\href{https://arxiv.org/abs/2609.10703v1}{arXiv:2609.10703v1}.

\bibitem{data}
F. Lander and S. Trettel,
\emph{Low vertex flat tori}, numerical data.
\href{https://github.com/FabianLander/low-vertex-flat-tori}{GitHub repository}.

\bibitem{lt}
F. Lazarus and F. Tallerie,
\emph{A universal triangulation for flat tori},
Discrete Comput. Geom. 71 (2024), 278--307.
\href{https://doi.org/10.1007/s00454-023-00598-7}{Journal article}.

\bibitem{quintanar}
T. Quintanar,
\emph{An explicit PL-embedding of the square flat torus into \(\R^3\)},
J. Comput. Geom. 11 (2020), 615--628.
\href{https://jocg.org/index.php/jocg/article/view/3114}{Journal article}.

\bibitem{schewe}
L. Schewe,
\emph{Nonrealizable minimal vertex triangulations of surfaces},
Discrete Comput. Geom. 43 (2010), 289--302.
\href{https://arxiv.org/abs/0801.2582}{arXiv:0801.2582}.

\bibitem{schwartz}
R. E. Schwartz,
\emph{Vertex-minimal paper tori},
preprint, 2025, version 9, 2026.
\href{https://arxiv.org/abs/2507.14998v9}{arXiv:2507.14998v9}.

\bibitem{sl}
T. Sulanke and F. H. Lutz,
\emph{Isomorphism-free lexicographic enumeration of triangulated
surfaces and 3-manifolds},
European J. Combin. 30 (2009), 1965--1979.
\href{https://arxiv.org/abs/math/0610022v3}{arXiv:math/0610022v3}.

\bibitem{tsuboi}
T. Tsuboi,
\emph{On origami embeddings of flat tori},
J. Math. Sci. Univ. Tokyo 31 (2024), 105--126.
\href{https://www.ms.u-tokyo.ac.jp/journal/31-2-1.pdf}{Journal article}.

\end{thebibliography}

\end{document}
